\documentclass[10pt,reqno]{amsart}
\usepackage[margin=0.7in]{geometry}

\usepackage{indentfirst,csquotes}



\usepackage{amssymb,amsthm,amsmath}
\usepackage{xcolor,paralist,hyperref,fancyhdr,etoolbox}
\newtheorem{theorem}{Theorem}[section]
\newtheorem{definition}{Definition}[section]
\newtheorem{lemma}{Lemma}[section]
\newtheorem{proposition}{Proposition}[section]
\newtheorem{corollary}{Corollary}[section]
\newtheorem{conjecture}{Conjecture}[section]


\hypersetup{ colorlinks=true, linkcolor=black, filecolor=black, urlcolor=black }
\def\proof{\noindent {\it Proof. $\, $}}
\def\endproof{\hfill $\Box$ \vskip 5 pt }

\usepackage{lipsum}

\begin{document}
\title{Erd\H{o}s Problem 99 with Finite Multiplicity} %%%%%%%%%%%%
\author[N Stewart]{Naveen Stewart}
\date{\today}
\email{npstewart18@gmail.com}
\maketitle

\let\thefootnote\relax

\begin{abstract}
Let $A \subset \mathbb{R}^2$ be a set of $n$ points with minimum distance equal to 1, chosen to minimize the diameter of $A$. In 1994, Paul Erd\H{o}s conjectured that for all sufficiently large $n$, there must be three points in $A$ which form an equilateral triangle of size 1. This problem later came to be known as Erd\H{o}s Problem 99. We prove that there must exist three points forming an equilateral triangle of size 1 if the number of total incongruent diameter-minimizing sets of $n$ points, denoted $h(n)$, is finite and $n$ is sufficiently large. This provides a partial solution to Erd\H{o}s Problem 99.
\end{abstract}

\section{Introduction}

Let:
\begin{equation}
    X \subseteq \mathbb{R}^2, \quad \text{sep}(X) = \min_{x \ne y} |x - y|, \quad D_n = \min \{\text{diam } X : |X| = n, \text{ sep}(X) \ge 1\}
\end{equation}
\begin{equation}
    \text{Opt}(n) = \{X : \text{diam } X = D_n, |X| = n, \text{sep} (X) \ge 1\}, \quad h(n) = \# (\text{Opt}(n) / \text{Isom}(\mathbb{R}^2))
\end{equation}

We prove:
\begin{theorem}
    There exists absolute lower bound $N$ such that if $n \ge N$ and $h(n) < \infty$, every $X \in \text{Opt}(n)$ contains three points at pairwise distance 1.
\end{theorem}
This solves Erd\H{o}s Problem 99 \cite{1} for all sufficiently large $n$ and finite $h(n)$.


\section{Proof of Theorem 1.1}



\subsection{Setup and Metric Estimates}


Suppose Theorem 1.1 does not hold. Then take a sequence of triangle-free optimal $X$ as $n \to \infty$, each with finite $h(n)$. Denote $D = D_n$, $H = \text{conv } X$, and $b = \# \text{vert } H$. All further asymptotics we introduce reference this sequence, and subsequences when specified by context. All fixed depths, component bounds, and positive length thresholds obey the limit $D \to \infty$. Put:
\begin{equation}
    a_\star = \sqrt{3}/2, \qquad c = \pi / (2\sqrt{3}), \qquad J(x, y) = (-y, x)
\end{equation}
The minimum diameter is achieved via compactness after fixing one point (similar unit distance problems with diameter minimizing constraints have well established this). A naive construction via disjoint half-unit disks gives $n \le (D + 1)^2$, and therefore $D \to \infty$.
\begin{lemma}
    (Metric control). We have:
    \begin{equation}n = cD^2 + O(D)
    \end{equation}
    \begin{equation}b \le \text{Per } H \le \pi D
    \end{equation}
    \begin{equation}0 \le \pi D^2 / 4 - |H| = O(D).
    \end{equation}
    
    For a suitable origin,
    \begin{equation}
        B(0, D / 2 - O(\sqrt{D})) \subset H \subset B(0, D/2 + O(\sqrt{D}))
    \end{equation}
    Every boundary support normal differs from the radial direction by $O(D^{-1/4})$. Hull sides with directions in $M$ intervals of width $a$ number $O(DM(a + D^{-1/4}))$. 
\end{lemma}
\begin{proof}
    Averaging translations of the unit triangular lattice over a disk of area $a_\star n$ provides $D^2 \le 4 a_\star n / \pi$. Via Oler's Inequality \cite{2}, $n \le |H| / a_\star + \text{Per } H / 2 + 1$, and isodiametry jointly give (4), (5), and (6). Bonnesen's inequality \cite{3,4}, $\pi^2 (R - r)^2 \le (\text{Per }H)^2 - 4 \pi |H| = O(D)$, gives $r \ge D / 2 - O(\sqrt{D})$.

    At an incenter, $|x| + r \le D$ for $x \in H$. A support normal $v$ at $x$ satisfies $x \cdot v \ge r$; this proves the angle bound. Normals in an interval therefore locate their boundary points in a radial arc of width $O(a + D^{-1/4})$; their projection onto its middle tangent provides a bound on its length by $O(D(a + D^{-1/4}))$. Each counted side has length at least 1. summing over the intervals gives $O(DM (a + D^{-1/4}))$ as the means of scaling.
\end{proof}


\subsection{Contact Geometry and Locking}



A feasible point motion curve preserves pairwise distances in $[1, D]$, and therefore does not alter the diameter of the set, producing optimizers. Among triangle-free set $X$, any allowed point motion contradicts finite $h(n)$ across the continuous motion curve; a distance squared from the set center and isometry invariance to this quantity proves such motions present an uncountably infinite number of incongruent sets. Note all candidates $X$ are noncollinear, as $D = O(\sqrt{n})$ whereas a collinear separated set has diameter at least $n - 1$.



\begin{lemma}
    (Contact geometry). Every boundary point is a strict hull vertex and a diameter endpoint. The unit-contact graph is connected and plane. Adding missing hull sides makes it 2-connected with strictly convex bounded faces. Every added side bounds one cap with a simple unit-edge return. Old bounded faces are convex unit polygons; quadrilateral faces are rhombi with smaller angle in $(\pi / 3, \pi / 2]$.
\end{lemma}

\begin{proof}
    We utilize the following release. For a proper nonempty label set $S$, suppose a unit vector $v$ has nonnegative scalar product with all crossing unit vectors and nonpositive product with all crossing diameter vectors, oriented into $S$. If all diameter products are strict, translate $S$ by $tv$. Otherwise, there is at most one tangent diameter vector, call it $De, e \perp v$: two different parallel diameter segments would make $|a + De|, |a - De| \le D$ for $a \ne 0$ impossible. Translate by $tv - \gamma t^2 e$ where $1 / (2D) < \gamma < 1/2$. The distance squared change at a crossing vector $a$ is:
    \begin{equation}
        2t, a \cdot v + t^2 (1 - 2 \gamma a \cdot e) + \gamma^2 t^4
    \end{equation}
    It has the required sign at tangent unit and diameter pairs. Other active signs are strict and inactive slacks persist. The sum of squared distances is a nonconstant polynomial. Either motion contradicts finite multiplicity.

    A boundary point without a diameter partner moves outward, contradicting $h(n) < \infty$. An interior point whose contact neighbors lie in a closed half-plane moves away from them. Interior contact fans have gaps less than $\pi$, as a result. A diameter endpoint is strictly exposed, as $|p_i - p_j| = D$ implies $(p_i - p_j) \cdot (p_k - p_i) \le - |p_k - p_i|^2 / 2 < 0$ where $k \ne i$. Unit edges can not cross, as two adjacent endpoints at a proper crossing have distances to the crossing summing to at most one, which breaks separation of points.

    A contact component missing the hull translates freely. Otherwise the components induce a noncrossing partition of the cyclic hull vertices. If there are several blocks, at least two are cyclic intervals (induct inside of the gaps of a block), and one has hull-normal arc of width at most $\pi$. This has no crossing contacts and its crossing diameter directions lie in that semicircle, so the release applies. The graph formed is connected. After adding the hull cycle, an articulation would remove a component with no other hull vertex. An extreme point of that component together with the articulation contradicts the interior fan condition. Face boundaries are simple and their angles less than $\pi$. Euler's formula and bounded incidences of added sides gives the exactly one added side per cap. A convex unit quadrilateral is a rhombus, and its shorter diagonal $2 \sin (\alpha / 2) > 1$ is in the triangle-free case.
\end{proof}

Let $m, t$ count unit and diameter pairs. Let $L$ denote the outer contact-walk length, with bridges counted twice. Let $q$ count rhombi, $f_k$ count the other bounded $k$-gonal faces, $f = \sum_{k\ge5}f_k$, and $\sigma = \sum_{k \ge 5} (k - 4) f_k$.

\begin{lemma}
    (Fan inertia and locking). For a complete simplicial fan in $\mathbb{R}^3$ with ray matrix $V$, a symmetric matrix $W$ with nonpositive off-diagonal entries supported on fan edges and $WV = 0$ has at most one negative eigenvalue. As a result, every feasibly first-order velocity of $X$ integrates into a feasible curve, where:
    \begin{equation}
        m = 2n - 2 - (L + \sigma) / 2, \qquad L + \sigma \le 2t \le 2b, \qquad q = n - O(D).
    \end{equation}
\end{lemma}

\begin{proof}
    Positive diagonal congruence normalizes the rays to $u_i$. Have $c_{ij} = W_{ij}$. Have gradient jumps $c_{ij}(u_i \times u_j)$ across oriented spherical edges. They close around vertices as $u_i \times \sum_j c_{ij}u_j = 0$; vertex cycles generate the dual cycle space of the sphere. The gradients integrate to a continuous convex homogeneous piecewise-linear function $h$, the support function of their convex hull $K$. The segment dual to edge $ij$ has length $c_{ij} |u_i \times u_j|$.

    Let $K_r = \{x: u_i \cdot x \le h(u_i) + \tau\}$. This is bounded, full dimensional and tends to $K$; every facet is present, as $w + \tau u_i$ strictly satisfies the other inequalities when $w$ is exposed at $u_i$. Minus its volume Hessian $W_\tau$ has at most one negative eigenvalue, annihilates the ray matrix, and has off-diagonal entries $-(\text{length}(F_i \cap F_j)) / |u_i \times u_j|$ for adjacent faces, based on \cite{5} in Theorem 2.4 Appendix A.3. These converge to $W_{ij}$. At an interior point of a positive original dual edge, every other active normal lies in its normal plane outside of the short arc $ij$. On compact interior subsegments the displacement $\tau (u_i + u_j) / (1 + u_i \cdot u_j)$ satisfies every other shifted inequality strictly, giving the lower limit. Compactness then gives the upper limit. Zero dual edges collapse. If a nonedge has a limiting exposed segment, its short spherical arc contains a fan ray $u_k = au_i + bu_j$ with $a, b > 0$ and $a + b > 1$. The shifted $i, j$ planes violate the shifted $k$ inequality. Other nonedges collapse. Opposite normals have disjoint facets. The kernel equations determine the diagonal limits. The inertia conclusion is closed under limits.

    Write $A_U, A_T$ for half distance squared derivative matrices and $(L(\omega)x)_i = \sum_j \omega_{ij} (x_i - x_j)$. For a proper stress $\lambda, \mu \ge 0$ with $A_U^T \lambda = A_T^T \mu$, lift the points to $v_i^\pm = \pm (p_i, 1)$. Triangulate both caps of their convex hull retaining the plane unit graph and all lifted points. A diameter $ij$ is an exposed cross-cap edge: the functional $(p_i - p_j) \cdot x - \frac{1}{2} (p_i - p_j) \cdot (p_i + p_j)z$ is maximized exactly at $v_i^+$ and $v_j^-$ Complete the triangulation and use a radial projection. With $C_\mu$ the diameter adjacency matrix and $B_0 = L(\lambda) - \text{diag} (C_\mu \textbf{1})$, the matrix:
    \begin{equation}
        W = \begin{pmatrix} B_0  & -C_\mu \\ -C_\mu & B_0 \end{pmatrix}
    \end{equation}
    has the fan support and annihilates its ray matrix. If $\mu \ne 0$, the symmetric subspace has the negative direction $(\textbf{1}, \textbf{1})$ with value $-8 \sum \mu$. The antisymmetric restriction is $2\Omega$, where $\Omega = L ( \lambda ) - L ( \mu )$. Fan inertia gives $\Omega \succeq 0$. For $\mu = 0$ this is the Laplacian inequality.

    On an affine slice transverse to the Euclidean motions, the active inequalities are $Az + Q(z) \ge 0$ with $A = 2(A_U ; - A_T)$ and $Q = (|z_i - z_j|^2_U ; -|z_i - z_j|^2_T)$. The stress sign implies $\eta \cdot Q(z) \ge 0$ if $\eta \ge 0, A^T \eta = 0$. Farkas' lemma puts $Q(z)$ in $\mathcal{C} = \{y : \exists w, Aw + y \ge0\}$. This polyhedral cone has a continuous piecewise-linear Lipschitz selection $S$ with $S(0) = 0$, $AS(y) + y \ge 0$: cut into coordinate orthants, triangulate compatibly, choose feasible corrections on rays, and interpolate. For $Av \ge 0$, the map $z \to tv + S(Q(Z))$ contracts a ball of radius $Rt^2$ about $tv$ for fixed sufficiently large $R$ and small $t$. Its continuous fixed point is feasible, with derivative $v$. Inactive slacks persist.

    Finite multiplicity therefore makes the velocity cone zero on the $(2n - 3)$-dimensional slice. Its $m + t$ signed active rows positively span that space, so $m + t \ge 2n - 2$. Euler and the incidence counts give $q + f = m - n + 1, 2m = L + 4(q + f) + \sigma$, producing the first identity in (9). The planar diameter bound $t \le b$ from \cite{6} Theorem 3 gives the second identity, while $f \le \sigma = O(D)$ gives the last.
    
\end{proof}



\subsection{Tracks and One Boundary Lemma}

For hull returns of $j$ unit edges, write $c_j$ for their number, $u = c_1$,
\begin{equation}
    \rho = \sum_{j \ge 3}(j - 2)c_j, \quad L = 2b - u + \rho, \quad \Delta = u - \sigma - \rho \ge 0.
\end{equation}
An old bounded unit $k$-gon has area at least $a_\star (k - 2) / 2$ by Oler. As $\sum_F (k_F - 2) = 2n - 2 - L$ and $|H| \le a_\star n$, the sum of all face-area excesses and cap areas is $O(D)$. In particular:
\begin{equation}
    q = cD^2 + O(D), \qquad \sum_{Q \text{ rhombus}} (\alpha_Q - \pi / 3) = O(D)
\end{equation}

\begin{lemma}
    (Tracks, clustering, and sources). Joining opposite edges across rhombi produces $s = O(D)$ tracks, each with two terminals on exceptional faces or the outer walk, and:
    \begin{equation}
        s = (L + \sigma + 4f) / 2 = b - \Delta / 2 + 2f
    \end{equation}
    A track's midpoint spine has unit steps in a cone of half-angle $\pi / 6$ about a fixed normal, with normal projection at least $a_\star$; its rhombus count $r(T) \le D / a_\star$. The crossing graph is simple and triangle-free. For cutoff $C$ and tolerance $\epsilon$, deleting $O(D / (\epsilon C))$ tracks laves the long-track directions in $O(s / C)$ intervals of width $O(\epsilon)$ modulo $\pi / 3$. For $C = D^{19/20}$, all but $o(D)$ unit and two-edge hull returns have a selected short track. These tracks are distinct and can not end on another return of length $o(D)$.
\end{lemma}



\begin{proof}
    Opposite rhombus edges are parallel; each spine step is another unit side with normal projections $\sin \alpha_Q \ge a_\star$. The projection keeps its sign, and as a result tracks neither close nor self-intersect, and $r a_\star \le D$. Terminal counting gives (13). Crossing directions have circular distance greater than $\pi / 3$ modulo $\pi$. For their normal angle $\gamma \in (\pi / 3, \pi / 2]$, the spines are graphs over the bisector with disjoint slope intervals: one ends at $\tan(\pi/6  - \gamma/2) < 0$, the other starts at $\tan(\gamma/2 - \pi/6) > 0$. They meet at most once. Three directions can not be pairwise more than $\pi / 3$ apart on a circle of length $\pi$. By (12), only $O(D / \epsilon)$ crossings have angle excess greater than $\epsilon$. Delete long tracks with at least $C / 2$ such incidences. Every survivor has at least $L / 2$ distinct good neighbors. Directions separated by more than $2 \epsilon$ moduli $\pi / 3$ have disjoint good-neighbor sets; a maximal separated family has size at most $2s / C$.

    Take $C = D^{19/20}, \epsilon = D^{-1/4}$. There are $O(D^{3/10}$ discarded tracks and $O(D^{1/20})$ remaining direction intervals. Metric control bounds unit sources on long tracks by $(D^{4/5})$. A two-edge return with two long ports has its chord direction in $O(D^{1/10})$ intervals of width $O(D^{-1/4})$, as the chord has length at least one. These return number $O(D^{17/20})$.

    If a selected track ends on a return of length $B_0 = o(D)$, its spine and paths through the terminal caps form a crosscut. The shore $S$ along the shorter hull arc is nonempty and proper, of diameter $O(C + B_0 + \sqrt{D}) = o(D)$, near radius $v$. All crossing unit vectors into $S$ equal one vector $z$. With $\tau = Jv$, the source gives $|z \cdot \tau| \ge 2 / 5$: each side of a two-edge return projects at least $1/2$ on its chord, and a unit return projects one. Every crossing diameter vector is $Dv + o(D)$. Translate $S$ with velocity:
    \begin{equation}
        V = \text{sgn}(z \cdot \tau) \tau - v / 4
    \end{equation}
    Then $z \cdot v \ge 2/5 - 1/4 = 3/20$, whereas each crossing diameter product is $-D / 4 + o(D) < 0$. Strict active signs yield a feasible nonconstant motion. This also excludes a track selected by two sources.
\end{proof}

Orient bounded-face ports counterclockwise. The two ports of a track on bounded faces have opposite vectors; a hull-source vector, oriented along the outer contact walk, equals its bounded-face terminal vector. These follow cancellation across rhombi.


\begin{lemma}
    (Boundary ports). Let a bounded-face port $z$ start a track at $x = (D / 2)v + o(D)$, and set $\tau = Jv, n = -J z$. Uniformly: (i) $z \cdot \tau \ge 1 / 2 - o(1)$ implies $r(T) = o(D)$; (ii) for fixed $\lambda > 0$, three ports on tracks of length at least $\lambda D$, within $o(D)$ of one boundary point, have sum of norm at least $n_\lambda > 0$; (iii) ports with $\tau(T) \ge \lambda D$ and $z \cdot \tau \ge - \epsilon$, where $0 < \epsilon < a_\star \lambda / 2$, belong to distinct tracks. For a family of $P$ such ports, after $\lambda, \epsilon \downarrow 0$ sufficiently slowly, their tracks touch at most:
    \begin{equation}
        Q_P \le \min \{DP / 2, cD^2 / 4\} + o ( D^2 )
    \end{equation}
    rhombi.
\end{lemma}

\begin{proof}
    For the other midpoint $y$, write $w = (y - x) / r = sn + t J n$. Then $s \ge a_\star, |t| \le 1 / 2, |w| \le 1$. Put $a = n \cdot v = z \cdot \tau = (J n) \cdot v$. Roundness gives:
    \begin{equation}
        r^2 |w|^2 + Dr (as + bt) \le o(D)^2.
    \end{equation}
    If $a \ge 1 / 2 - o(1)$, then $as + bt \ge - o(1)$, since $a a_\star - \sqrt{1 - a^2} / 2 \ge 0$ for $a \ge 1 / 2$. As $|w| \ge a_\star$, this proves (i). Three unit vectors summing to zero are $2 \pi / 3$ apart, and one has tangent projection at least $1 / 2$; compactness and (i) prove (ii).

    For (iii), opposite selected ends would have $n \cdot v_x \ge - \epsilon, n \cdot v_y \le \epsilon$, so $a_\star \lambda D \le n (y - x) \le \epsilon D + o(D)$ impossible. Also:
    \begin{equation}
        \frac{|t|}{s^2 + t^2} \le \frac{1}{2}, \qquad \frac{s}{s^2 + t^2} \le \frac{1}{a_\star}
    \end{equation}
    As a result, (16) gives $r \le ( 1 / 2 + \epsilon / a_\star) D + o(D)$ when $a \ge - \epsilon$. Every unit direction $u$ of a forward spine displacement lies within $\pi / 6$ of $n$, therefore $u \cdot v \ge -(1/2 + \epsilon) = - \mu_\epsilon$. Further $|(D / 2)v + tu|^2 \ge D^2 (1 - \mu_\epsilon^2) / 4$ for $t \ge 0$. All touched rhombi lie in this common annulus, up to $o(D)$ in radius; they have disjoint interiors and area at least $a_\star$. Therefore $Q_P \le c \mu_\epsilon^2 D^2 + o(D^2)$ and $Q_P \le (1 / 2 + \epsilon / a_\star) DP + o(D^2)$. Since $P \le s = O(D)$, slow passage $\epsilon \downarrow 0$ gives (15).
\end{proof}


\subsection{Propagation and the Short-track Count}


Write $p = f_5, h = f_6$ and $Q_9 = \sum_{k \ge 9}(k - 8) f_k$. All exceptional-face and hull-return incidences total $O(D)$.

\begin{lemma}
    (Local capacities). In a convex unit polygon with unit-separated vertices, mark the edges directed within $\delta = o(1)$ of one vector. The maximum marked counts at $k = 5, 6, 7, 8$ are $2, 2, 3, 3$; for $k \ge 9$ the count is at most $(k - 4) - (k - 8)/4$. With two such pentagon sides, each other direction differs by $O(\delta)$ from $2 \pi / 3$, $\pi$, or $4 \pi / 3$ relative to the marked direction. With three such heptagon edges, two others are within $O(\sqrt{\delta})$ of the opposite direction. A $j$-edge return has no edges near the negative hull tangent for $j \le 3$ and at most $2(j - 2) / 3$ for $j \ge 4$.
\end{lemma}

\begin{proof}
    For $r$ marked polygon edges, projection of closure gives $r(1 + \cos \delta) \le k$. Three marked hexagon edges or four octagon edges force a limiting straight polygon with repeated vertices, forbidden by separation. For $k \ge 9$, use $\cos \delta \ge 9 / 10$ and $k / 1.9 \le (3k - 8) / 4$. Convex cyclic order makes the marked edges consecutive in the extremal cases. For a pentagon write its consecutive marked vertices as $A = ( - \cos \beta, 0), B = (0, - \sin \beta), C = (\cos B, 0), 2 \beta \le \delta$, and the others as $(x, y)$, $(z, w)$. The unit-edge and separation conditions give $\cos \beta, x + \sin \beta, y \ge \cos (2\beta)/ 2 $ and $-cos \beta, z + \sin \beta, w \le \cos (2 \beta) / 2$, Their mutual distance one forces $x = 1/2 + O (\beta), z = -1/2 + O(\beta)$. Circle equations and convexity give $y, w = \sqrt3 / 2 + O(\beta)$. For a heptagon, the four consecutive marked vertices are $(i, 0) + O(\delta), 0 \le i \le 3$. The remaining vertices $X_1, X_2, X_3$ satisfy $(X_1)_x) \le 5 / 2 - O(\delta)$ and $(X_3)_x \le 1/2 + O(\delta)$ by separation from adjacent marked vertices. Both intervening unit steps have the horizontal projection $-1 + O(\delta)$. For a return with $r$ negative projections of magnitude at least $d$, its chord gives $1 \le j - r - dr$. This excludes $j \le 2$. For $j = 3$ the limiting vectors would be $e, e, -e$ ordered in some way, repeating a vertex. Taking $d \ge 9 / 10$ gives the integer bounds one and two for $j = 4, 5$. For $j \ge 6$, we have $(j - 1) / 1.9 \le 2(j-2) / 3$.
\end{proof}


\begin{proposition}
    (Boundary count). At one cutoff $C = o(D)$, outside of $o(D)$ face incidences, pentagons have at least three short ports and hexagons at least four. Marked short tracks give each such face a distinct retained unit-hull source label at distance $o(D)$. The marked vectors sum to $\tau + o(1)$ at pentagons and $2 \tau + o(1)$ at hexagons. As activated pentagon has three marked ports and all five ports short; other pentagons have one marked port, hexagons two.

    Let $E_C$ count short tracks, $\beta$ the cycle rank in hull-containing components of their terminal graph, $R = s - E_C$, $A$ the number of activated pentagons, and $C_1 \le c_2$ the two-edge returns with one short port. After deleting marked tracks, good faces have degrees two to four, with total surplus $Y$ above degree two. Then:
    \begin{equation}
        \Delta, \rho, \sum_{k\ge 7} k f_k = o(D), \quad u = p + 2h + o(D), \quad f = p + h + o(D), \quad b = p + 2h + c_2 + o(D),
    \end{equation}
    \begin{equation}
        E_C = b + f + \beta + o(D), \quad R = f - \beta + o(D), \quad f + c_2 \le \pi D + o(D),
    \end{equation}
    \begin{equation}
        Y = 2 (\beta - A) + C_1 + o(D), \quad C_2 \le 2c_2 - C_1 + o(D),
    \end{equation}
    where $C_2$ counts unused short face-to-cap edges.
\end{proposition}





\begin{proof}
    \textit{Stability}. Send distinct labels along the $u - o(D)$ short unit-hull source tracks. A pentagon retains its first arrival and emits the second on its almost opposite edge; a heptagon retains two and emits the third on either opposite edge. Each emits at most once, retaining the last arriving label on its output; other receivers stop tokens. Set $e_j = (20 \cdot 4^j)^{-1}, \Lambda _ j = D^{1 - e_j}$. In round $j$ keep only output tracks shorter than $\Lambda_j$. Track cancellation makes incoming face vectors near the positive source tangent and outputs near its negative; arrivals on outer returns point near the negative hull tangent. For fixed depth $K$, locality gives distance $O_K (\Lambda_{j-1})$, input error $O_K (D^{-e_{j-1}})$ and output error $O_K (D^{-2e_j})$. No kept output reuses a source or output track: its endpoint vectors would be opposite while both near the same negative tangent. Labels never split.

    Clustering at $\Lambda_j$, tolerance $D^{-1/4}$, omits $O(D^{1/4 + e_j})$ long tracks. The other long directions form $O(D^{e_j})$ families; distinct labels, the output error and hull spread bound long output attempts by $O_K (D^{1-e_j} + D^{3/4 + e_j}) = o_K (D)$. Each omitted track has two terminals. For large receivers put $a = D^{-e_K / 4}$, $B_K = D^{1-e_K / 8}$. Three source positions separated circularly by at least $a$ would force receiver area at least $c_o D^2 a^3 - O_K (D\Lambda_K + D^{3/2}) \gg D$, which contradicts the area budget. Therefore, its sources lie in at most two arcs of width $O(a)$, and there are $O_K (D^{1 - e_K / 4})$ of them. There are $O(D / B_K)$ receivers with more than $B_K$ edges. Their total arrivals are $o_K (D)$. These are discarded. At every other receiver source tangents agree to $o_K (1)$.

    Local capacities bound cap absorption by $2 \rho / 3 + o_K (D)$, face absorption by $\sigma - f_7 - f_8 - Q_9 / 4 + o_K (D)$. If $F_j$ counts kept outputs, conservation gives:
    \begin{equation}
        F_j \ge \Delta + \rho / 3 + f_7 + f_8 + Q_9 / 4 - o_K (D), \qquad \sum_{j=1}^K F_j \le p + f_7.
    \end{equation}

    Divide by $K$, let $D \to \infty$, then $K \to \infty$, and use $k \le 9 (k - 8)$ for $k \ge 9$. This proves (18).

    \textit{Saturation}. Rerun with only pentagons emitting. At a round with $F_j \le p / K$, conservation and $u = \sigma + o(D)$ leave at most $p / K + o_K (D)$ unused absorption slots. Almost every pentagon receives one input and every hexagon two. Keep the first source label at each face permanently. These labels are injective, and all such faces are in an $o(D)$ boundary layer. An activated pentagon has two positive inputs and one negative output, resulting in marked sum $\tau + o(1)$. Fix $\lambda > 0$. Outside of $o(D)$ terminals, macroscopic port directions $(r \ge \lambda D)$ lie in finitely many shrinking intervals. Put $\zeta = e^{i \theta} = \tau$ and let $Z_j$ sum $j$ such vectors. A pentagon with only one or two short ports satisfies $\zeta + Z_4 = o(1)$ or $| \zeta + Z_3| = 1 + o(1)$. Lemma 2.5 (ii) bounds $|Z_3|$ away from zero, so these restrict $\theta$ to finitely many shrinking intervals. For a hexagon with two or three short ports use $2 \zeta + Z_4 = o(1)$ or $|2\zeta + Z_3| = 1 + o(1)$, the second of which requires $|Z_3| \ge 1 - o(1)$. Injective labels and hull spread make these cases $o(D)$. At an activated pentagon every remaining direction is fixed relative to $\theta$ by Lemma 2.6. A long port again pins $\theta$. Slow choices $K \to \infty$ with $\lambda \downarrow 0$ yield a common sublinear cutoff dominating marked tracks and boundary errors.

    \textit{Connection}. Terminal-graph nodes are faces and hull returns, tracks its edges. Orient caps oppositely to the outer walk. Cancellation over any node set gives the sum of outgoing port vectors as minus the sum of its hull-side vectors. Fix $M$ and $0 < e \le 1 / 20$, and put $L_0 = D^{1 - e}, e' = e / (12M), L_1 = D^{1 - e'}$. Discard $L_0$-components of at most $M$ nodes that meet high faces ore return with at least three edges. This costs $o_M (D)$ nodes. Remaining nodes have at most six ports. A hull-containing such component with not track in $[L_0, L_1)$ has outgoing sum close to most of $O_M (D^{6Me'})$ vectors by clustering, apart from $o_M (D)$ exceptions. Its hull-side directions differ by $O_M (D^{-e} + D^{-1/4})$ and their sum has norm at least $1/2$. Closure pins these directions; spread bounds their number by $O_M (D^{1 + 6Me' - e} + D^{3/4 + 6Me'}) = o_{M, e} (D)$. Repeat $J$ times with $e_{j+1} = e_j / (12M)$. If $r_j$ counts hull nodes in components of at most $M$ nodes, it decreases and each new edge serves at most $2M$ of them, therefore $Jr_J \le 2MS + o_{M, J}(D)$. There are at most $(b + f) / (M + 1)$ larger components. Choose $M$, then $J$, then a slow diagonal, where there are $o(D)$ hull components. Good faces connect to their source by marked tracks, so only $o(D)$ other face ports and components remain. Euler's graph identity gives $E_C = b + f + \beta + o(D)$. We have (13) and (18) together give (19) fully.

    \textit{Degrees}. Let $X_+$ be the short-degree excess over three at pentagons and four at hexagons. Short-degree summation gives $2E_C = u + 2c_2 - C_1 + 3p + 4h + X_+ + o(D)$, and therefore $X_+ = 2 \beta + C_1 + o(D)$. Ordinary pentagons lose one marked incidence, hexagons lose two, and activated pentagons lose three. There are at most $2c_2 - C_1 + o(D)$ other unused cap ports, which proves (20). Enlarging the cutoff while keeping it sublinear preserves these statements while deleting at most $sC = o(D^2)$ crossings overall.
\end{proof}


\subsection{Local Algebra and the Cycle Inequality}



Enlarge $C = o(D)$ such that intermediate tracks are negligible:
\begin{equation}
    \lim_{\lambda \downarrow 0} \lim_{D \to \infty} \sup D^{-1} \# \{T : C \le r(T) < \lambda D\} = 0
\end{equation}
The $D^{-1} \sum_T \delta_{r(T) / D}$ has bounded mass and compact support. Pass to a weak limit $\mu$ and choose continuity points $a_j \downarrow 0$. A slow diagonal $C / D = a_{j(D)}$, dominating earlier cutoffs, captures mass $\mu (\{0\})$ below $C / D$. At a fixed continuity point $\lambda$, the intermediate mass tends to $\mu([0, \lambda)) - \mu(\{0\})$, which vanishes as $\lambda \downarrow 0$. Fix $\lambda$ and write $N_\lambda = \# \{T : C \le r(T) < \lambda D\}$. Discard faces meeting intermediate tracks, at a cost $O(N_\lambda) + o(D)$ ports. All equations below carry this error, absorbed by (22). Take $D \to \infty$ first, component bounds to infinity next, and $\lambda \downarrow 0$ last.

Delete marked tracks from the short terminal graph. Let $\mathcal{H}$ be good faces of degree two and $\mathcal{G}$ be those of degree three or four. An ordinary face in $\mathcal{H}$ has two retained (nonshort) ports, an activated one none; degree-three faces have one and degree-four faces none. Marked sums are $k \tau + o(1)$ with $k = 1, 2$. Identify the plane with $\mathbb{C}$, writing $\zeta = e^{i \theta} = \tau$. In a bounded short-track pattern tangents agree to $o(1)$. If a port from a neighbor into an $\mathcal{H}$ face is $z$ in the neighbor's orientation, and the receiver's retained-vector sum is $Z$, its other unused port is $z - a$, where $a = k \zeta + Z$. Unit constraints imply:
\begin{equation}
    z^2 - az + a / a^\star = 0, \qquad a = k / \zeta + \bar{Z}.
\end{equation}
Here $\star$ is conjugation on the unit circle, and then algebraic continuation. At infinity its branches are $z = k \zeta + O(1)$ or $z = 1 / \bar{Z} + o(1)$ if $Z \ne 0$. Otherwise for $Z = 0$, $z = d \zeta$ with $d^2 - kd + 1 = 0$. All leading coefficients have nonnegative real part. Let $\mathcal{J} \subset \mathcal{G}$ be degree-three faces incident with $\mathcal{H}$. Choose an $\mathcal{H}$ response $z$ and let $v$ be the retained unit vector. Either other incoming port satisfies (22) with $a_J = \ell \zeta + v + z$, $a_J^\star = \ell / \zeta + \bar{v} + 1 / z$. Every branch has $a_J = A \zeta + O(1), \text{Re }A > 0$. The limiting constant of $a_J^\star$ is $\bar{v}$ or $\bar{v} + \bar{Z}$, nonzero by Lemma 2.5 (ii). These responses again have leading coefficient zero or $A$.

\textit{Counting principle}. In a bounded pattern, clustering places macroscopic vectors near a fixed finite list after a subsequence. If a necessary equation is not an identity on any branch of its finitely many circle quadratics, multiplication over root choices and clearing denominators gives a nonzero polynomial in $\zeta$ of bounded degree. Denominator zeros are finite too. Compactness puts approximate solutions in a bounded number of shrinking neighborhoods of these roots, provided the noncancellation margins persist. Injective source labels and hull spread then count only $o(D)$ such patterns. Expansions at infinity test identities; they are not physical motions.

\begin{lemma}
    (Algebraic exclusions). Fix $0 < \epsilon < a_\star \lambda / 2$. Flag a face if a retained port has $z \cdot \tau \ge - \epsilon$ and let $P$ count all these ports. Up to the specified errors, there are no $\mathcal{H}-\mathcal{H}$ incidences; an unflagged $\mathcal{J}$ face has at most one unflagged $\mathcal{H}$ incidence; adjacent unflagged $\mathcal{J}$ faces can not have an unflagged $\mathcal{H}$ incidence. For $U = \mathcal{G}$, $\mathcal{K} := \mathcal{G} \setminus \mathcal{J}$, let $e_U, b_U$ count internal short edges and short edges to caps. Then $e_U + b_U \ge |U| - o(D)$. If $h_C$ counts $\mathcal{H}-cap$ incidences, $h_C \le 2 (c_2 - C_1) + o(D)$.
\end{lemma}





\begin{proof}
    For adjacent $\mathcal{H}$ faces, put $a = k \zeta + Z, b = \ell \zeta + W$. Their three unit vectors require:
    \begin{equation}
        aa^\star bb^\star (a + b) (a^\star + b^\star) + (ab^\star - a^\star b)^2 = 0.
    \end{equation}

    The $\zeta^3$ coefficient is $k \ell (k + \ell) \bar{ZW} (\bar{Z} + \bar{W})$. If $Z = 0 \ne W$, the $\zeta^2$ coefficient is $k^2 [\ell (k + \ell) + 1]\bar{W}^2 \ne 0$, symmetrically for $W = 0 \ne Z$. If both vanish, the expression becomes $k^2 \ell^2 (k + \ell)^2 > 0$. If $W = -Z \ne 0$, put $r = |Z|, x = \text{Re}(\zeta \bar{Z})$. Division by $(k + \ell)^2$ gives:
    \begin{equation}
        (k^2 + r)^2 (\ell^2 + r^2) - 4r^2 + 2 \{k (\ell^2 + r^2) - \ell (k^2 + r^2)\} x + 4 (1 - k \ell) x^2.
    \end{equation}


    It vanishes identically only for $k = \ell = 1, r = 1$. The retained pairs then form opposite pairs with gaps $\pi / 3, 2 \pi /3$. One has tangent projection at least $1 / 2$, prohibited by Lemma 2.5 (i). A loop gives $k \zeta + Z = 0$. The counting principle excludes the rest.

    Two unflagged $\mathcal{H}$ responses $z_1, z_2$ at an unflagged $\mathcal{J}$ face require $a a^\star = 1$, with $a = \ell \zeta + v + z_1 + z_2 = A \zeta + O(1), \text{Re }A > 0$. The constant $B$ in $a^\star = \ell / \zeta + \bar{v} + 1 / z_1 + 1 / z_2$ is $\bar{v}$ plus zero, one or two conjugate retained-pair sums. All have strictly negative common tangential projection before conjugation, so $B \ne 0$ and $a a^\star$ can not be identically one. For adjacent such $\mathcal{J}$ faces, choose one unflagged $\mathcal{H}$ response at each. In (24), $a_i = A_i \zeta + O(1), a_i^\star = B_i + o(1)$, with $\text{Re } A_i > 0$ and $B_1 B_2 (B_1 + B_2) \ne 0$. The first term has nonzero cubic coefficient, the squared term order at most two. A $\mathcal{J}$ loop reduces to $| \ell \zeta + v| = 1$. Repeated incidences obey the same necessary equations.

    For a bounded component of $\mathcal{G}$ without a cap edge, sum closure: outgoing $\mathcal{H}$ responses satisfy $\sum_i x_i = - k_0 \zeta - Z_0, k_0 > 0$. Their leading coefficients have nonnegative real part, a contradiction. For a component of $\mathcal{K}$, neighbors lie in $\mathcal{H}$ or $\mathcal{J}$. Use a chosen $\mathcal{H}$ response at every such $\mathcal{J}$ and the two-stage expansion. Component of each fixed bounded size number $o(D)$. The larger ones number $O(D / m)$. There are at most $b_U$ cap-touching components. Therefore, $e_U \ge |U| - \# \text{components}(U) \ge |U| - b_U - o(D)$. Components touching discarded faces are charged to discarded ports, with an absolute constant independent of $m$.

    Finally, a $C_1$ cap feeding $\mathcal{H}$ has a macroscopic other port $v$ in the outer-walk orientation. Reflection in its chord gives short-port vector $\zeta^2 \bar{v} + o(1)$. Receiver closure requires $|k \zeta + \zeta^2 \bar{z} + Z|^2 - 1 = o(1)$. Its constant Fourier coefficient is $k^2 + |Z|^2 > 0$. The counting principle excludes these incidences, and every remaining two-edge cap contributes at most two.
\end{proof}



\begin{proposition}
    (Joint count). For the flagged ports just chosen:
    \begin{equation}
        \beta + \frac{2}{5}P + \frac{2}{5} c_2 + \frac{1}{5}C_1 \ge \frac{2}{5} f + A - o(D).
    \end{equation}
\end{proposition}


\begin{proof}
    Suppress accounted errors. Let $g_1, g_2$ count degree-three and degree-four faces. Then $Y = g_1 + 2g_2, f = |\mathcal{H}| + Y - g_2$. Degree summation and $e_\mathcal{G} + b_\mathcal{G} \ge |\mathcal{G}|$ gives:
    \begin{equation}
        2 | \mathcal{H} | \le Y + C_2, \qquad g_2 \le \frac{3}{2} Y + \frac{1}{2} C_2 - f.
    \end{equation}
    Let $j_i$ count $\mathcal{J}$ faces with $i$ incidences to $\mathcal{H}$, let $x \le 4 g_2$ count $\mathcal{H}-$-degree-four edges, and let $e_{JJ}$ count $\mathcal{J}$-$\mathcal{J}$ edges. The degree equations and the connection for $\mathcal{K}$ give:
    \begin{equation}
        2 | \mathcal{H}| = j_1 + 2j_2 + 3j_3 + x + h_C, \qquad 3 j_1 + 2 j_2 + j_3 + x \le Y + 2 e_{JJ} + C_2 - h_C.
    \end{equation}
    Therefore, $M = e_{JJ} + 2j_2 + 3j_3 \ge 3f - 7/2 Y + 3 g_2 - x - h_C - C_2$. Assign each $\mathcal{J}$-$\mathcal{J}$ edge to an endpoint that is flagged or has all $\mathcal{H}$ incidences flagged; the exclusions must guarantee one. A $\mathcal{J}$ face with $i$ such incidences also receives charge $2(i - 1)$. A flagged $\mathcal{J}$ receives at most $(3 - i) + 2(i - 1) = i + 1 \le 4$. An unflagged $\mathcal{J}$ has at most one unflagged $\mathcal{H}$ incidence. If it has one, it receives no edge charge and passes charge two per flagged incidence. Otherwise distribute its charge over all $i$ flagged incidences, at most $(i + 1) / i \le 2$ each. Each flagged $\mathcal{H}$ receives at most four units through its two incidences. Every charged flagged face has a flagged port, so $M \le 4P + o(D)$. Therefore:
    \begin{equation}
        4f \le 5Y + 4P + h_C + C_2 + o(D) \le 10(\beta - A) + 4P + 4c_2 + 2C_1 + o(D),
    \end{equation}
    which is therefore (26). Choose continuity points $\lambda_j \downarrow 0$ with intermediate mass below $1 / j$, then $0 < \epsilon_j < a_\star \lambda_j / 2$, and component bounds tending to infinity. Increase successive $D$ thresholds so fixed-parameter errors are below $D / j$ and $L / D < \lambda_j$. This gives one common family of ports for (26) and (15).
\end{proof}





\subsection{The Resulting Contradiction}

From (19), (26), $C_1 \le c_2$, and $A \ge 0$, we have:
\begin{equation}
    R - P \le \frac{3}{5} (f - P) + \frac{2}{5}c_2 + \frac{1}{5}C_1 - A + o(D) \le \frac{3}{5} (\pi D - P) + o(D).
\end{equation}
Delete the short tracks, losing $o(D^2)$ crossings, and then the $P$ selected tracks. The remaining crossing graph is simple and triangle-free. Mantel's bound $e \le v^2 / 4$ follows by summing $\deg x + \deg y \le v$ over edges and using Cauchy: $4e^2 / v \le \sum \deg^2 \le ve$. The annular half of (15) gives that:
\begin{equation}
    \frac{3}{4} cD^2 - o(D^2) \le (R - P)^2 / 4, \qquad R - P \ge \sqrt{3c} D - o(D)
\end{equation}
Along a subsequence $P / D \to x$, (30) implies that $0 \le x \le a := \pi - \frac{5}{3} \sqrt{3c}$. The other half of (15) gives:
\begin{equation}
    c \le F(x) := \frac{9}{100} (\pi - x)^2 + x / 2
\end{equation}
But $F$ is convex, and both endpoint values on $[0, a]$ are below $c$. With $157 / 50 < \pi < 22 / 7$ and $\sqrt{3} < 87/50$ gives $c > 9 / 10$, and $F(0) < \frac{9}{100}(22 / 7)^2 < 9 / 10 < c$. Also $\sqrt{3c} > 41 / 25$ so $a < 22 / 7 - 41 / 15 = 43 / 105 < 9 /20 < c / 2$, and $F(a) = 3c / 4 + 1 / 2 < c$. This yields a contradiction. If $a < 0$, the earlier bounds already contradict each other. The positive endpoint margin is approximately $0.0186353$.

Therefore, no unbounded sequence of triangle-free minimizers with finite multiplicity exists. This proves Theorem 1.1.



\section{Lean 4 Formalization}

The previous work has been formalized in Lean 4 under standard axioms and with no proof gaps or assumptions. The Lean 4 proof is available at \url{https://github.com/Naveen18888/Erdos-99-with-Finite-Multiplicity}.


\section{Future Work}

This completes the proof of Erd\H{o}s Problem 99 for finite $h(n)$. The $h(n) = \infty$ case is unresolved.

\begin{thebibliography}{99}

\bibitem{1}
T.~Bloom
\newblock Erd\H{o}s Problem \#99.
\newblock \url{https://erdosproblems.com/99}, 2023.
\newblock [Online]
\bibitem{2} Oler N. A Finite Packing Problem. Canadian Mathematical Bulletin. 1961;4(2):153-155. doi:10.4153/CMB-1961-018-7
\bibitem{3} Bonnesen, T. Über eine Verschärfung der isoperimetrischen Ungleichheit des Kreises in der Ebene und auf der Kugeloberfläche nebst einer Anwendung auf eine Minkowskische Ungleichheit für konvexe Körper. Math. Ann. 84, 216–227 (1921). https://doi.org/10.1007/BF01459405
\bibitem{4} Zhang, Z., Zhou, J. Bonnesen-style Wulff isoperimetric inequality. J Inequal Appl 2017, 42 (2017). https://doi.org/10.1186/s13660-017-1305-3
\bibitem{5} Izmestiev, I. The Colin de Verdi\`ere number and graphs of polytopes. arXiv. 2008. https://arxiv.org/abs/0704.0349
\bibitem{6} Erd\H{o}s P. On sets of distance of $n$ points, Amer. Math. Monthly \textbf{53}(5) (1946), 248-250

\end{thebibliography}
\end{document} 